%% kapa: exam=1 round=31 year=2020 period=2 form=A subject=appraisal_law subject_ko=감정평가관계법규 seq=5 number=5
%% source: Q-Net 감정평가사 기출문제 (cst00309) artlSeq=5207219 file=extracted/1008/5207219/A형/A형/2020년도 제31회 감정평가사 1차 시험_2교시_A형.pdf
%% conversion: kordoc 4.24.51 -> tools/gen1.py (auto); review=proofread-sonnet; flags=none
\begin{question}{5}
국토의 계획 및 이용에 관한 법령상 입지규제최소구역 지정에 관한 설명으로 옳은 것을 모두 고른 것은?
\begin{framedtext}
ㄱ. 지역의 거점 역할을 수행하는 철도역사를 중심으로 주변지역을 집중적으로 정비할 필요가 있는 지역은 입지규제최소구역으로 지정될 수 있다. \par
ㄴ. 세 개 이상의 노선이 교차하는 대중교통 결절지로부터 3킬로미터에 위치한 지역은 입지규제최소구역으로 지정될 수 있다. \par
ㄷ. 「도시 및 주거환경정비법」상 노후․불량 건축물이 밀집한 공업지역으로 정비가 시급한 지역은 입지규제최소구역으로 지정될 수 있다. \par
ㄹ. 입지규제최소구역계획에는 건축물의 건폐율․용적률․높이에 관한 사항이 포함되어야 한다. \par
\end{framedtext}
\begin{choices}
  \item ㄱ, ㄴ
  \item ㄱ, ㄷ
  \item ㄴ, ㄹ
  \item ㄱ, ㄷ, ㄹ
  \item ㄴ, ㄷ, ㄹ
\end{choices}
\end{question}
\answer{4}
